\section{部分公式的证明}
\subsection{静电场的旋度与环路定理}
关注公式 (\ref{eq:2.1 E v}) 和 (\ref{eq:1.4 nabla 2 r r'}), 对于任意体积元 $v$, $\nabla\dfrac{1}{|\bm{r}-\bm{r'}|}$ 仅与 $\bm{r}$ 有关, 因此微分算子可以移出被积函数之外:
\begin{equation}
    \bm{E}(\bm{r})=-\frac{1}{4\pi\epsilon_0}\int_V\rho(\bm{r'})\nabla\frac{1}{|\bm{r}-\bm{r'}|}\mathrm{d}v=-\nabla\left[\frac{1}{4\pi\epsilon_0}\int_V\frac{\rho(\bm{r'})}{|\bm{r}-\bm{r'}|}\mathrm{d}v\right].
\end{equation}

上式两边取旋度, 得
\begin{equation}
    \nabla\times\bm{E}(\bm{r})=-\nabla\times\nabla\left[\frac{1}{4\pi\epsilon_0}\int_V\frac{\rho(\bm{r'})}{|\bm{r}-\bm{r'}|}\mathrm{d}v\right].
\end{equation}

显然 $\nabla\left[\dfrac{1}{4\pi\epsilon_0}\displaystyle\int_V\dfrac{\rho(\bm{r'})}{|\bm{r}-\bm{r'}|}\mathrm{d}v\right]$ 是梯度, 结果为标量. 标量的旋度应为 0. 所以
\begin{equation}
    \nabla\times\bm{E}=0.
\end{equation}
此为公式 (\ref{eq:2.1 rot}).

上式两边对任意曲面 $\bm{S}$ 求积分:
\begin{equation}
    \int_S\nabla\times\bm{E}\cdot\mathrm{d}\bm{S}=0.
\end{equation}

根据斯托克斯定理 (\ref{eq:1.4 stokes}),
\begin{equation}
    \int_S\nabla\times\bm{E}\cdot\mathrm{d}\bm{S}=\oint_C\bm{E}\cdot\mathrm{d}l.
\end{equation}
所以,
\begin{equation}
    \oint_C\bm{E}\cdot\mathrm{d}l=0.
\end{equation}
此为公式 (\ref{eq:2.1 rot int}).

\subsection{电磁场能流密度与能量密度} \label{proofs: 2}
认为 $f$ 即为洛伦兹力和库仑力的合力, 也即
\begin{equation}
    \begin{aligned}
        \bm{f}\cdot\bm{v} & =\rho(\bm{E}+\bm{v}\times\bm{B})\cdot\bm{v}                                                                                  \\
                          & =\bm{E}\cdot\rho\bm{v}                                                                                                       \\
                          & =\bm{E}\cdot\bm{J}                                                                                                           \\
                          & =\bm{E}\cdot\nabla\times\bm{H}-\bm{E}\cdot\frac{\partial\bm{D}}{\partial t}                                                  \\
                          & =-\nabla\cdot(\bm{E}\times\bm{H})+\bm{H}\cdot\nabla\times\bm{E}-\bm{E}\cdot\frac{\partial\bm{D}}{\partial t}                 \\
                          & =-\nabla\cdot(\bm{E}\times\bm{H})-\bm{H}\cdot\frac{\partial\bm{B}}{\partial t}-\bm{E}\cdot\frac{\partial\bm{D}}{\partial t}.
    \end{aligned}
\end{equation}

代入式 (\ref{eq:2.8 S}), 可以对比得到
\begin{gather}
    \bm{S}=\bm{E}\times\bm{H}, \\
    \frac{\partial w}{\partial t}=\bm{E}\cdot\frac{\partial\bm{D}}{\partial t}+\bm{H}\cdot\frac{\partial\bm{B}}{\partial t}.
\end{gather}
利用麦克斯韦方程组, 可以进一步推得
\begin{equation}
    w=\frac{1}{2}(\bm{E}\cdot\bm{D}+\bm{H}\cdot\bm{B}).
\end{equation}
